\documentclass[10pt,a4paper]{amsart}
\usepackage[margin=19mm]{geometry}
\usepackage{amsmath,amssymb,mathtools}
\usepackage[scr=rsfs,cal=euler]{mathalfa}
\usepackage{xfrac}
\usepackage[unicode,hidelinks,bookmarks=false]{hyperref}
\newcommand{\ie}{i.e.\ }
\newcommand{\cf}{cf.\ }
\newcommand{\resp}{respectively\ }
\newcommand{\Subsection}{\subsection}
\providecommand{\nobreakdash}{\nobreak\hbox{-}}
\setlength{\emergencystretch}{2em}
\multlinegap0pt
\usepackage{amssymb}
\usepackage{enumerate}
\newtheorem{theorem}{Theorem}
\newtheorem{lemma}{Lemma}
\newtheorem{corollary}{Corollary}
\usepackage{cleveref}
\crefname{theorem}{Theorem}{Theorems}
\crefname{lemma}{Lemma}{Lemmas}
\crefname{corollary}{Corollary}{Corollarys}
\newcommand{\R}{\mathbb{R}}
\renewcommand{\d}{\textnormal{d}}
\title{The concentration-compactness principle for the nonlocal anisotropic $p$-Laplacian of mixed order}
\author{Jamil Chaker, Minhyun Kim, Marvin Weidner}
\date{Source: arXiv:2105.13612v2 (CC BY 4.0)}
\begin{document}
\maketitle

\noindent\emph{Source and licence.} The concentration-compactness principle for the nonlocal anisotropic $p$-Laplacian of mixed order. Version arXiv:2105.13612v2 (CC BY 4.0), distributed by arXiv under the Creative Commons Attribution 4.0 International licence, http://creativecommons.org/licenses/by/4.0/, as recorded in the arXiv licence metadata for this exact version and shown on the abstract page https://arxiv.org/abs/2105.13612v2. Full licence notice: \texttt{licenses/arxiv-2105.13612-CC-BY-4.0.txt}.\par\smallskip

\noindent\textit{Editorial scope.} Counted unit: the article's theorem thm:unit\_norm (source0.tex lines 437--440). The article's complete original proof of it is delivered with it (source0.tex lines 774--825, one \texttt{proof} environment of 52 lines, which the article places away from the statement). No mathematics is written, repaired or re-derived: the statement and the proof are the article's own lines, in the article's own notation, and no retained line is substituted or re-flowed.  Retained uncounted as the source-local context the proof needs: lines 421--423; lines 429--431; lines 610--613; lines 617--622; lines 628--631; lines 689--695; lines 720--731.  Nothing was omitted from any retained span, so the package carries no omission record.  No retained line contains a citation, so the wrapper carries no bibliography and no bibliography entry.  No retained line contains a figure environment, an image-inclusion command or drawing code, so no image is delivered and the package declares no asset dependency.  Editorial formatting only, in addition: an AMS \texttt{amsart} preamble that restates the article's own declarations and macros used by the retained lines, namely the environments \texttt{theorem}, \texttt{lemma}, \texttt{corollary} on separate counters, together with the standard packages those lines need.  The retained environments print in this document as Theorem 1, Lemma 1, Corollary 1 in order of appearance, whereas the article prints the same items with the numbering of its own full document.  The complete native source of the article as distributed in the arXiv e-print for this exact version is bundled unchanged as \texttt{sources/111967/source0.tex}.
\begin{equation} \label{eq:sup}
\int_{M_1(0)} |v_k|^{p^{\ast}} \,\mathrm{d}x = \frac{1}{2} = \sup_{y \in \mathbb{R}^n} \int_{M_1(y)} |v_k|^{p^{\ast}} \,\mathrm{d}x.
\end{equation}

\begin{equation} \label{eq:conv_ptwise}
v_k \to v ~ \text{a.e. in}~ \mathbb{R}^n
\end{equation}

\begin{lemma}\label{lem:inftyrevHolder} 
For $\mu_{\infty}$ and $\nu_{\infty}$ as above, the following is true:
\[S\nu_{\infty}^{p_{\max}/p^{\ast}} \leq \mu_{\infty}.\]
\end{lemma}

\begin{align}
\label{eq:nuChar}
\nu_{\infty} &= \lim_{R \to \infty} \limsup_{k \to \infty} \int_{\mathbb{R}^n} |w_k|^{p^{\ast}} \psi_R^{p^{\ast}} \,\mathrm{d}x \quad\text{and}\\
\label{eq:muChar}
\mu_{\infty} &= \lim_{R \to \infty} \limsup_{k \to \infty} \sum_{i=1}^n \frac{1}{p_i} \int_{\mathbb{R}^n} |D^{s_i}_{p_i}w_k|^{p_i} \psi_R^{p_i} \,\mathrm{d}x.
\end{align}

\begin{equation}
\label{eq:nuinftyrep1}
\nu_{\infty} = \lim_{R \to \infty} \limsup_{k \to \infty} \int_{\mathbb{R}^n} |v_k|^{p^{\ast}} \psi_R^{p^{\ast}} \,\mathrm{d}x.
\end{equation}

\begin{corollary}
\label{cor:CC}
There exists an at most countable set $\Xi \subset \R^n$ and positive weights $(\nu_{\xi})_{\xi \in \Xi}$, $(\mu_{\xi})_{\xi \in \Xi}$ such that
\begin{equation*}
\nu = \sum_{\xi \in \Xi} \nu_{\xi}\delta_{\{\xi\}} \quad\text{and}\quad \mu \ge \sum_{\xi \in \Xi} \mu_{\xi}\delta_{\{\xi\}}.
\end{equation*}
\end{corollary}

\begin{corollary}
\label{cor:CC2}
Every $\xi \in \Xi$ is an atomic point of $\widetilde{\mu}$ and it holds
\begin{equation}
\label{eq:finallemmaHelp2}
S \nu_{\xi}^{\frac{p_{\max}}{p^{\ast}}} \le \widetilde{\mu}(\{\xi\}).
\end{equation}
In particular,
\begin{equation*}
\Vert \widetilde{\mu} \Vert \ge S \left(\sum_{\xi \in \Xi} \nu_{\xi}^{\frac{p_{\max}}{p^{\ast}}} +  \Vert v \Vert_{L^{p^{\ast}}(\R^n)}^{p_{\max}}\right).
\end{equation*}
\end{corollary}

\begin{theorem} \label{thm:unit_norm}
Let $v:\R^n\to\R$ be the defined as above. Then
\[\|v\|_{L^{p^{\ast}}(\mathbb{R}^n)} = 1.\]
\end{theorem}

\begin{proof} [Proof of \Cref{thm:unit_norm}]
Let us assume by contradiction that $\Vert v \Vert_{L^{p^{\ast}}(\R^n)} < 1$.
First, we have that
\begin{equation}\label{eq:unit_normHelp1}
\begin{aligned}
1 = \frac{1}{S} \left(\Vert \widetilde{\mu}\Vert + \mu_{\infty}\right) & \ge \Vert v \Vert_{L^{p^{\ast}}(\R^n)}^{p_{\max}} + \sum_{\xi\in\Xi} \nu_{\xi}^{\frac{p_{\max}}{p^{\ast}}} + \nu_{\infty}^{\frac{p_{\max}}{p^{\ast}}} \\
& \ge \Vert v \Vert_{L^{p^{\ast}}(\R^n)}^{p_{\max}} + \Vert \nu \Vert^{\frac{p_{\max}}{p^{\ast}}} + \nu_{\infty}^{\frac{p_{\max}}{p^{\ast}}},
\end{aligned}
\end{equation}
where the second step follows from \Cref{lem:inftyrevHolder} and \Cref{cor:CC2}. In order to prove the first equality in \eqref{eq:unit_normHelp1} we define $\psi_R : \R^n \to \R$ as in the proof of \Cref{lem:inftyrevHolder} and compute
\begin{equation*}
\begin{split}
S &= \lim_{k \to \infty} \sum_{i=1}^n\frac{1}{p_i} \Vert D_{p_i}^{s_i}v_k \Vert_{L^{p_i}(\R^n)}^{p_i}\\
&= \lim_{R \to \infty} \lim_{k \to \infty} \int_{\R^n} \sum_{i=1}^n \frac{1}{p_i}\vert D^{s_i}_{p_i}v_k\vert^{p_i}(1-\psi_R^{p_i}) \,  \d x + \lim_{R \to \infty} \limsup_{k \to \infty} \int_{\R^n} \sum_{i=1}^n \frac{1}{p_i}\vert D^{s_i}_{p_i}v_k\vert^{p_i} \psi_R^{p_i} \, \d x\\
&= \lim_{R \to \infty} \int_{\R^n} (1-\psi_R^{p_i}) \, \d \widetilde{\mu} + \lim_{R \to \infty} \limsup_{k \to \infty} \int_{\R^n} \sum_{i=1}^n \frac{1}{p_i}\vert D^{s_i}_{p_i}w_k\vert^{p_i} \psi_R^{p_i} \, \d x\\
&= \Vert \widetilde{\mu} \Vert + \mu_{\infty},
\end{split}
\end{equation*}
where we used \eqref{eq:muChar}. This proves \eqref{eq:unit_normHelp1}.
On the other hand, it holds that
\begin{equation}
\label{eq:unit_normHelp2}
\Vert v \Vert_{L^{p^{\ast}}(\R^n)}^{p_{\max}} + \Vert \nu \Vert^{\frac{p_{\max}}{p^{\ast}}} + \nu_{\infty}^{\frac{p_{\max}}{p^{\ast}}} \ge 1.
\end{equation}
This follows from the observation
\begin{equation}
\label{eq:unit_normHelp3}
\begin{split}
1 &= \lim_{k \to \infty} \Vert v_k \Vert_{L^{p^{\ast}}(\R^n)}^{p^{\ast}} = \lim_{R \to \infty} \lim_{k \to \infty} \int_{\R^n} (1-\psi_R^{p^{\ast}}) \vert v_k \vert^{p^{\ast}} \, \d x + \lim_{R \to \infty} \lim_{k \to \infty} \int_{\R^n} \psi_R^{p^{\ast}} \vert v_k \vert^{p^{\ast}} \, \d x\\
&=\Vert v \Vert_{L^{p^{\ast}}(\R^n)}^{p^{\ast}} + \Vert \nu \Vert+ \nu_{\infty},
\end{split}
\end{equation}
where we used that $\vert v_k \vert^{p^{\ast}}\, \d x \rightharpoonup \vert v \vert^{p^{\ast}} \, \d x + \nu$ as a consequence of Brezis–Lieb's lemma and \eqref{eq:conv_ptwise}, and \eqref{eq:nuinftyrep1}. In order to derive \eqref{eq:unit_normHelp2}, we take \eqref{eq:unit_normHelp3} to the power $p_{\max}/p^{\ast} < 1$.
Combining \eqref{eq:unit_normHelp1} and \eqref{eq:unit_normHelp2} yields that
\begin{equation*}
\Vert v \Vert_{L^{p^{\ast}}(\R^n)}^{p_{\max}} + \Vert \nu \Vert^{\frac{p_{\max}}{p^{\ast}}} + \nu_{\infty}^{\frac{p_{\max}}{p^{\ast}}} = 1.
\end{equation*}
Consequently, $\Vert v \Vert_{L^{p^{\ast}}(\R^n)}^{p^{\ast}}, \Vert \nu \Vert, \nu_{\infty} \in \{0,1\}$. By assumption, we can now deduce that $v = 0$. Moreover, by \eqref{eq:sup} we know that $\nu_{\infty} \le 1/2$ and therefore it must be that 
\begin{equation*}
1 = \Vert \nu \Vert = \sum_{\xi \in \Xi} \nu_{\xi},
\end{equation*}
where we applied \Cref{cor:CC} in the last step.
From \eqref{eq:unit_normHelp1}, we deduce that $1 \ge  \sum_{\xi\in\Xi} \nu_{\xi}^{\frac{p_{\max}}{p^{\ast}}}$. It follows that
\begin{equation*}
1 = \left( \sum_{\xi \in \Xi} \nu_{\xi} \right)^{\frac{p_{\max}}{p^{\ast}}} \le \sum_{\xi \in \Xi} \nu_{\xi}^{\frac{p_{\max}}{p^{\ast}}} \le 1,
\end{equation*}
and therefore all inequalities in the above line can be replaced by equality signs. Thus, $\nu_{\xi} \in \{0,1\}$ for every $\xi \in \Xi$ and therefore, there exists exactly one $\xi_0 \in \Xi$ such that $\nu_{\xi_0} = 1$ and $\nu_{\xi} = 0$ for every $\xi \in \Xi \setminus \{\xi_0\}$. This leads to a contradiction since by \eqref{eq:sup}:
\begin{equation*}
\frac{1}{2} = \sup_{y \in \R^n} \int_{M_1(y)} \vert v_k \vert^{p^{\ast}} \, \d x \ge \sup_{y \in \R^n} \int_{M_1(\xi_0)} \vert v_k \vert^{p^{\ast}} \, \d x \to \Vert \nu \Vert = 1.
\end{equation*}
Therefore, it must hold that $\Vert v \Vert_{L^{p^{\ast}}(\R^n)} = 1$, as desired.
\end{proof}

\end{document}
